% Intended LaTeX compiler: pdflatex
\documentclass[a4paper,12pt]{article}
\usepackage[utf8]{inputenc}
\usepackage[T1]{fontenc}
\usepackage{graphicx}
\usepackage{longtable}
\usepackage{wrapfig}
\usepackage{rotating}
\usepackage[normalem]{ulem}
\usepackage{amsmath}
\usepackage{amssymb}
\usepackage{capt-of}
\usepackage{hyperref}
\usepackage{\string~/.emacs.d/common}
\author{mahmood}
\date{\today}
\title{comparison sort}
\hypersetup{
 pdfauthor={mahmood},
 pdftitle={comparison sort},
 pdfkeywords={},
 pdfsubject={},
 pdfcreator={Emacs 28.2 (Org mode 9.5.5)}, 
 pdflang={English}}
\begin{document}

\maketitle
\begin{definition}
a \textbf{comparison sort} is a type of \textbf{sorting algorithm} that only reads the list elements through a single abstract comparison operation (often a "less than or equal to" operator or a three-way comparison) that determines which of two elements should occur first in the final sorted list.\\
\begin{lemma}
any comparison-sort based algorithm has a \href{20221130014441-time_complexity.org}{time complexity} that is bound by \(\Omega(n\log n)\)\\
\begin{my_example}
a vector \texttt{A} of size \texttt{n} is \textbf{defined} to be a "hill" iff for every \(i \le i < j \le m\) we know \(A[i] \le A[j]\) and for every \(m \le i < j \le n\) we know \(A[i] \geq A[j]\) where \(m\) is a constant and \(i,j\) are variable indicies\\
prove that any comparison-based algorithm that turns a vector into a hill has to perform \(\Omega(n\log n)\) operations\\
\begin{answer}
we assume in contradiction that there exists a comparison-based algorithm that turns a vector into a hill whose time complexity is less than \(\Omega(n\log n)\)\\
consider the following function:\\
\includesvg{/home/mahmooz/brain/out/30t764s}\\

each line from 3 to 6 runs in \(\Theta(n)\) time, except the first which we know by the assumption runs in less than \(\Theta(n\log n)\), so the whole function runs in less than \(\Theta(n\log n)\) time\\
this is a comparison-based function that sorts a given array in less than \(\Theta(n\log n)\) time, which contradicts with our theorem that no sorting function can have a lower time complexity than \(\Theta(n\log n)\)\\
we arrived at a contradiction and so it has to be true that this function is bound from below by \(\Theta(n\log n)\), i.e. its time complexity is \(\Omega(n\log n)\)\\
\end{answer}
\end{my_example}
\label{orgde41dd1}
\end{lemma}
\end{definition}
\end{document}